% ---------------------------------------------------------------------------
%  QML Lab research project 12: QCNN for Phase Classification
%  Build:  latexmk -pdf main.tex   (needs qmlproject.sty, harvardqml.bst, references.bib)
% ---------------------------------------------------------------------------
\documentclass[11pt,a4paper]{article}
\usepackage{qmlproject}
\qmlheader{Project 12: QCNN for Phase Classification}

\hypersetup{pdftitle={QML Lab project 12: QCNN for Phase Classification}}

\begin{document}

\begin{center}
  {\small\textcolor{black!55}{QML Lab \textbullet\ research project 12}}\\[0.5em]
  {\LARGE\bfseries QCNN for Phase Classification}\\[0.6em]
  {\large Recognising the phases of the transverse-field Ising model directly from quantum states.}\\[0.6em]
  {\small Living document \textbullet\ last updated \today}
\end{center}

\begin{projectcard}
\small
\begin{tabularx}{\linewidth}{@{}>{\bfseries}l X >{\bfseries}l X@{}}
Quadrant  & \QQ\ (quantum data, quantum processor) & Track & Quantum data \\
Difficulty & Hard & Tools & PennyLane, SciPy \\
Builds on & Modules 5 and 8 of the QML Lab & Workspace & \path{qml_projects/qcnn-phase-classification/} \\
\end{tabularx}
\end{projectcard}

\paragraph{Summary.}
A quantum convolutional neural network (QCNN) processes a many-body quantum state with layers of
local unitaries and pooling, much like a convolutional network processes an image. Its input is the
quantum state itself, which makes this the lab's main \QQ\ project. The QCNN learns to tell the
ordered from the disordered phase of the transverse-field Ising chain, and is compared with classical
learners that only see measurement snapshots.

% ===========================================================================
\section{Background and key concepts}
% ===========================================================================
\subsection{The transverse-field Ising model}
The chain of $N$ spins with Hamiltonian
\begin{equation}
  H=-J\sum_{i=1}^{N-1}Z_iZ_{i+1}-h\sum_{i=1}^{N}X_i
\end{equation}
has a quantum phase transition at $h=J$ in the thermodynamic limit \citep{pfeuty1970,sachdev2011}:
for $h<J$ the ground state is ferromagnetically ordered, for $h>J$ it is paramagnetic. In finite
chains the transition is smoothed and shifted, which a classifier must cope with.

\subsection{Quantum convolutional neural networks}
\citet{cong2019} introduced the QCNN. A \emph{convolution} layer applies the same two-qubit unitary to
neighbouring pairs; a \emph{pooling} layer measures (or traces out) half of the qubits and applies
conditioned rotations to their neighbours; after $O(\log N)$ rounds a final unitary acts on the
remaining qubit, whose $\langle Z\rangle$ gives the phase. The network has only $O(\log N)$ parameters.
\citet{pesah2021} proved that QCNNs do not suffer from barren plateaus, and \citet{herrmann2022}
ran a QCNN on a superconducting processor to recognise a symmetry-protected phase.

\subsection{Quantum versus classical learners}
A classical learner, as in \citet{carrasquilla2017}, must work from measurement snapshots, losing
phase information at every measurement. The QCNN acts on the state before measurement. Comparing both
with the same number of state copies is the core of the \QQ\ question.

% ===========================================================================
\section{Research questions}
% ===========================================================================
The project is open-ended: these questions are the starting point, not the finish line. Add new
ones to the README and the journal as they come up.

\begin{rqlist}
  \item How close to the critical point $h=J$ does the QCNN classify correctly, and how does that change with $N$?
  \item Do QCNN gradients stay large as the chain grows, as theory predicts?
  \item How many state copies does the QCNN need compared with a classical model trained on measurement snapshots?
\end{rqlist}

% ===========================================================================
\section{Method}
% ===========================================================================
\subsection{Data}
Ground states of the Ising chain for $N=4$ and 8 (later 12) computed by exact diagonalisation with
SciPy sparse methods, for $h/J$ from 0.1 to 2.0. Training uses states away from the transition
($h<0.8$ and $h>1.2$); testing covers the full range.

\subsection{Models}
QCNN with parameterised two-qubit convolution blocks, pooling by controlled rotations and a final
single-qubit readout; binary cross-entropy on the phase label; Adam. Input states are loaded exactly
with \texttt{StatePrep} in simulation.

\subsection{Baselines}
Logistic regression on two-point correlations $\langle Z_iZ_j\rangle$ (an order-parameter-like
feature); an MLP on computational-basis snapshots, as in \citet{carrasquilla2017}; an unstructured
hardware-efficient classifier with the same parameter count.

\subsection{Metrics}
\begin{itemize}
  \item Accuracy far from and near the transition; predicted transition point $h^*$ where $p=0.5$.
  \item Gradient variance versus $N$.
  \item Number of state copies or shots to reach a given accuracy.
\end{itemize}

% ===========================================================================
\section{Planned experiments}
% ===========================================================================
Each experiment becomes one \texttt{[[experiments]]} entry in the workspace, so its runs, parameters
and metrics are stored and compared in the portal.

\begin{table}[h]
\centering\small
\begin{tabularx}{\linewidth}{@{}l>{\raggedright}X>{\raggedright}X>{\raggedright\arraybackslash}X@{}}
\toprule
ID & Experiment & What varies & Main metrics \\
\midrule
  E1 & Small chain & $N=4$ & accuracy, $h^*$ \\
  E2 & Larger chain & $N=8$ (then 12) & accuracy, $h^*$ \\
  E3 & Near criticality & test $h\in[0.8,1.2]$ & accuracy \\
  E4 & Copies & shots per state & accuracy vs snapshot MLP \\
  E5 & Trainability & QCNN vs unstructured ansatz & gradient variance \\
\bottomrule
\end{tabularx}
\caption{Experiment plan. Results go into the journal and the results log (\cref{sec:log}).}
\end{table}

% ===========================================================================
\section{Running it in the QML Lab}
% ===========================================================================
The workspace lives in \path{qml_projects/qcnn-phase-classification/}. Scripts import three helpers from
\path{qml_lab/plots.py}:
\begin{itemize}
  \item \texttt{params()} reads the parameters from the portal's run form (defaults when run by hand);
  \item \texttt{save()} stores figures with the run;
  \item \texttt{log\_metrics()} records numbers that appear in the run table.
\end{itemize}

\begin{lstlisting}
qml_projects/qcnn-phase-classification/
  project.toml          metadata + experiments
  README.md             goal, open questions, ideas
  experiments/          runnable scripts
  docs/main.tex         this document
\end{lstlisting}

First experiment: This experiment is \emph{planned}: write the script, add the block below to \texttt{project.toml}, then run \texttt{python manage.py sync\_projects}.

\begin{lstlisting}[language={}]
[[experiments]]
slug = "ising-phases"
title = "QCNN on transverse-field Ising ground states"
script = "experiments/e01_ising_phases.py"
figure = "phase_prediction.png"

[experiments.params]
n_spins = 8
h_min = 0.1
h_max = 2.0
n_train = 40
epochs = 100
seed = 0
\end{lstlisting}

\begin{lstlisting}[language=bash]
# from the portal folder: pick up changes to project.toml
python manage.py sync_projects
cd qml_projects/qcnn-phase-classification
# run by hand with the default parameters
PYTHONPATH=../../qml_lab python experiments/e01_ising_phases.py
\end{lstlisting}

% ===========================================================================
\section{Quality bar}
% ===========================================================================
\begin{itemize}
  \item A classical baseline with a similar number of trainable parameters, tuned with the same care \citep{bowles2024}.
  \item At least five random seeds per configuration, reported as mean $\pm$ standard deviation.
  \item How results change with the number of qubits.
  \item At least one run under a realistic noise model.
  \item Compare quantum and classical learners at an equal number of state copies.
\end{itemize}

% ===========================================================================
\section{Open directions}
% ===========================================================================
\begin{itemize}
  \item The symmetry-protected topological phase of the cluster-Ising model, the original QCNN task.
  \item Inputs compressed by the quantum autoencoder (Project 06).
  \item Classical shadows as the classical learner's input (Project 13).
\end{itemize}

% ===========================================================================
\section{Risks and limitations}
% ===========================================================================
\begin{itemize}
  \item In simulation the exact state is free; the cost of preparing it on hardware is the real bottleneck.
  \item Finite-size effects shift $h^*$ away from 1; this is physics, not a bug.
  \item Training sets are small; regularise and use several seeds.
\end{itemize}

% ===========================================================================
\section{Results log}
\label{sec:log}
% ===========================================================================
Summarise findings here as they become stable; day-to-day notes belong in the portal journal.

\begin{table}[h]
\centering\small
\begin{tabularx}{\linewidth}{@{}l l >{\raggedright}X >{\raggedright\arraybackslash}X@{}}
\toprule
Date & Run & Finding & Next step \\
\midrule
 & & & \\
 & & & \\
 & & & \\
\bottomrule
\end{tabularx}
\end{table}

\bibliographystyle{harvardqml}
\bibliography{references}

\end{document}
